\chapter{Lower bound (Section 4 and Appendix A)}
\label{chap:lower_bound}

\section{A lower bound for non-stationary bandits}

A non-stationary bandit with finitely many arms $\mathcal{A}$ is a sequence
$\nu = (\nu_{t, a})_{t \in \N, a \in \mathcal{A}}$ of probability measures on $\R$, played through
LML's environment \texttt{Environment.banditSeq}. For an identification algorithm with budget
$T$, $X_t$ is the arm of round $t$ and $\hat k$ the output.

\begin{lemma}[Non-stationary Kaufmann--Capp\'e--Garivier bound, Lemma 2]
  \label{lem:gen_lower}
  \lean{MaynardZhang2026Complexity.le_max_prob_ne}
  \leanok
  \uses{lem:bretagnolle_huber, lem:data_processing_output, lem:chain_rule_nonstationary}
  Let $\nu, \nu'$ be non-stationary bandits with $\KL(\nu_{t, a} \,\|\, \nu'_{t, a}) < \infty$ for
  all $t, a$, let $k \ne k'$ be arms and let $A$ be an identification algorithm with budget $T$.
  For every run of $A$ against $\nu$ (law $\Prob_\nu$) and every run against $\nu'$,
  \[
    \max\left(\Prob_\nu(\hat k \ne k), \Prob_{\nu'}(\hat k \ne k')\right)
    \ge \frac14 \exp\left(-\sum_{a} \sum_{t < T} \Prob_\nu(X_t = a)\,
      \KL(\nu_{t, a} \,\|\, \nu'_{t, a})\right).
  \]
\end{lemma}

\begin{proof}
  \leanok
  \uses{lem:bretagnolle_huber, lem:data_processing_output, lem:chain_rule_nonstationary}
  Let $E = \{\hat k \ne k\}$; since $k \ne k'$, $E^c \subseteq \{\hat k \ne k'\}$. By the
  Bretagnolle--Huber inequality (Lemma~\ref{lem:bretagnolle_huber}) for the laws of the outputs,
  $\Prob_\nu(E) + \Prob_{\nu'}(E^c) \ge \frac12 \exp(-\KL)$, so the maximum is at least
  $\frac14 \exp(-\KL)$, where $\KL$ is the divergence between the laws of the outputs. It is at
  most the divergence between the laws of the histories of the $T$ rounds
  (Lemma~\ref{lem:data_processing_output}), which is the sum in the statement
  (Lemma~\ref{lem:chain_rule_nonstationary}).
\end{proof}

\begin{remark}
  The paper states Lemma~\ref{lem:gen_lower} for the best arms $k$ and $k'$ of $\nu$ and $\nu'$,
  assumed different; only $k \ne k'$ is used. Its proof goes through the change of measure at a
  stopping time (Lemma~\ref{lem:change_of_measure}) and Lemma 19 of
  \cite{kaufmann2016complexity}; the route above, through the chain rule for the divergence of the
  histories, avoids both.
\end{remark}

\begin{lemma}[Change of measure at a stopping time, Lemma 6]
  \label{lem:change_of_measure}
  \lean{MaynardZhang2026Complexity.trajMeasure_eq_lintegral_exp_neg_llrProcess,
    Bandits.llrProcess,
    Bandits.trajMeasure_banditSeq_apply_eq_setLIntegral_of_isStoppingTime,
    Bandits.trajMeasure_banditSeq_apply_eq_setLIntegral,
    Bandits.feedback_banditSeq_eq_withDensity,
    Learning.IsAlgEnvSeq.hasLaw_history_withDensity_of_feedback_eq,
    Learning.stepKernel_eq_withDensity_of_feedback_eq}
  \leanok
  Let $\nu, \nu'$ be non-stationary bandits with countably many arms and $\nu_{t,a}$, $\nu'_{t,a}$
  mutually absolutely continuous, and let $\Prob_\nu$, $\Prob_{\nu'}$ be the laws of the
  trajectory of an algorithm against them. Let
  $L_n = \sum_{s < n} \log \frac{d\nu_{s, X_s}}{d\nu'_{s, X_s}}(Y_s)$ be the log-likelihood ratio
  of the first $n$ rounds, and $\mathcal{F}_n$ the $\sigma$-algebra generated by the rounds
  $0, \dots, n$. For every stopping time $\sigma$ of this filtration, finite everywhere, and every
  $E \in \mathcal{F}_\sigma$, $\Prob_{\nu'}(E) = \E_\nu[\indic_E \exp(-L_{\sigma + 1})]$.
\end{lemma}

\begin{proof}
  \leanok
  The step kernel of round $n$ against $\nu'$ has density
  $(h, (a, y)) \mapsto \frac{d\nu'_{n, a}}{d\nu_{n, a}}(y)
  = \exp(-\log \frac{d\nu_{n, a}}{d\nu'_{n, a}}(y))$ ($\nu_{n, a}$-almost everywhere, by mutual
  absolute continuity) with respect to that against $\nu$, so by induction on $n$ the law of the
  history of the first $n$ rounds against $\nu'$ has density $\exp(-L_n)$ with respect to that
  against $\nu$. This is the identity for $E \in \mathcal{F}_n$ (determined by the first $n + 1$
  rounds) with $L_{n + 1}$. For $E \in \mathcal{F}_\sigma$, $E$ is the disjoint union of the sets
  $E \cap \{\sigma = n\} \in \mathcal{F}_n$, on which $L_{\sigma + 1} = L_{n + 1}$.
\end{proof}

\begin{remark}
  \label{rem:change_of_measure_index}
  In the paper the rounds are numbered from $1$, and $\mathcal{F}_t$ and $L_t$ cover the same
  rounds $1, \dots, t$: its $L_\sigma$ is $L_{\sigma + 1}$ here. The statement written in phase 1
  of this formalization, with $L_\sigma$, was false (an indexing error of the translation, not of
  the paper): for the constant stopping time $\sigma = 0$ and the event
  $E = \{Y_0 \in B\} \in \mathcal{F}_0$, its left-hand side is $\E[\nu'_{0, X_0}(B)]$ and its
  right-hand side $\Prob_\nu(E) = \E[\nu_{0, X_0}(B)]$.
\end{remark}

\section{The optimization problem}

\begin{definition}[Optimization problem of Lemma 3]
  \label{def:opt_value}
  \lean{MaynardZhang2026Complexity.PairFeasible, MaynardZhang2026Complexity.pairValue,
    MaynardZhang2026Complexity.optValue}
  \leanok
  \uses{def:adjacency, def:design}
  For a pair $(x, x')$ of vertices, $(\theta, v)$ is \emph{feasible} if
  $(x - y)^\top \theta \ge \Delta$ for every vertex $y \ne x$ and $(x' - y)^\top (\theta + v) \ge
  \Delta$ for every vertex $y \ne x'$. For a matrix $A$, $p_A(x, x') = \inf\{v^\top A v : (\theta,
  v) \text{ feasible}\}$, and
  \[
    f(\X, \Delta) = \max_{\lambda \in \triangle_{\X}} \min_{(x, x') \in \Y} p_{A(\lambda)}(x, x'),
  \]
  the maximum being over all designs (a supremum in Lean).
\end{definition}

\begin{lemma}[Feasible points exist]
  \label{lem:pair_feasible_nonempty}
  \lean{MaynardZhang2026Complexity.exists_pairFeasible,
    MaynardZhang2026Complexity.exists_pairFeasible_of_finite,
    MaynardZhang2026Complexity.exists_forall_le_inner_sub_of_mem_vertices,
    MaynardZhang2026Complexity.pairValue_nonneg}
  \leanok
  \uses{def:opt_value, lem:strict_expose}
  For distinct vertices $x, x'$ there is a feasible $(\theta, v)$. For $A$ positive semidefinite,
  $p_A(x, x') \ge 0$.
\end{lemma}

\begin{proof}
  \leanok
  \uses{lem:strict_expose}
  Let $w$, $w'$ strictly expose $x$ and $x'$ (Lemma~\ref{lem:strict_expose}); for $c$ large,
  $\theta = c w$ and $\theta + v = c w'$ satisfy the finitely many constraints. The values
  $v^\top A v$ are nonnegative.
\end{proof}

\begin{lemma}[Cauchy--Schwarz inequality for Mahalanobis norms]
  \label{lem:mahalanobis_cauchy_schwarz}
  \lean{Matrix.PosDef.inner_sq_le_mahalanobisSq_inv_mul}
  \leanok
  \uses{def:design}
  For $A$ positive definite and $z, v \in \R^d$, $(z^\top v)^2 \le \norm{z}^2_{A^{-1}}
  \norm{v}^2_A$.
\end{lemma}

\begin{proof}
  \leanok
  $z^\top v = (A^{-1/2} z)^\top (A^{1/2} v)$, and $\norm{A^{-1/2} z}^2 = \norm{z}^2_{A^{-1}}$,
  $\norm{A^{1/2} v}^2 = \norm{v}^2_A$; apply the Cauchy--Schwarz inequality of $\R^d$.
\end{proof}

\begin{lemma}[Lower bound on the inner problem, Lemma 7]
  \label{lem:pair_value_lower}
  \lean{MaynardZhang2026Complexity.le_pairValue,
    MaynardZhang2026Complexity.PairFeasible.two_mul_le_inner,
    MaynardZhang2026Complexity.PairFeasible.div_mahalanobisSq_inv_le}
  \leanok
  \uses{def:opt_value, lem:pair_feasible_nonempty, lem:mahalanobis_cauchy_schwarz}
  Let $\Delta > 0$, $A$ positive definite and $(x, x') \in \Y$. Then
  $p_A(x, x') \ge 4\Delta^2 / \norm{x - x'}^2_{A^{-1}}$.
\end{lemma}

\begin{proof}
  \leanok
  \uses{lem:pair_feasible_nonempty, lem:mahalanobis_cauchy_schwarz}
  The constraints at $y = x'$ and at $y = x$ add up to $(x' - x)^\top v \ge 2\Delta$ for every
  feasible $(\theta, v)$. By the Cauchy--Schwarz inequality for the inner products of $A^{-1}$ and
  $A$ (Lemma~\ref{lem:mahalanobis_cauchy_schwarz}),
  $4 \Delta^2 \le ((x' - x)^\top v)^2 \le \norm{x - x'}^2_{A^{-1}} \norm{v}^2_A$. The infimum is
  over a nonempty set (Lemma~\ref{lem:pair_feasible_nonempty}).
\end{proof}

\begin{lemma}[Upper bound on the inner problem, Lemma 8]
  \label{lem:pair_value_upper}
  \lean{MaynardZhang2026Complexity.pairValue_le,
    MaynardZhang2026Complexity.exists_pairFeasible_mahalanobisSq_eq}
  \leanok
  \uses{def:opt_value, lem:is_adjacent_iff, lem:pair_feasible_nonempty}
  Let $A$ be positive definite and $(x, x') \in \I$. Then
  $p_A(x, x') \le 4\Delta^2 / \norm{x - x'}^2_{A^{-1}}$.
\end{lemma}

\begin{proof}
  \leanok
  \uses{lem:is_adjacent_iff, lem:pair_feasible_nonempty}
  Let $w$ and $\eps$ be given by Lemma~\ref{lem:is_adjacent_iff}, $c = \norm{x - x'}^2_{A^{-1}} > 0$,
  $\theta = \frac{\Delta}{c} A^{-1}(x - x') + \alpha w$ and $v = -\frac{2\Delta}{c} A^{-1}(x - x')$.
  Then $v^\top A v = 4\Delta^2 / c$, the constraints at $y = x'$ and $y = x$ hold with equality
  ($(x - x')^\top w = 0$), and the other constraints hold for $\alpha$ large since
  $(x - y)^\top w = (x' - y)^\top w \ge \eps$. The values are nonnegative, so the infimum is at most
  $4\Delta^2 / c$.
\end{proof}

\begin{remark}
  The paper assumes $\Delta > 0$ in Lemma~\ref{lem:pair_value_upper}; the construction works for
  every $\Delta$. The same construction with $\theta = -v/2 + \alpha w$ shows that, for an adjacent
  pair, the feasible $v$ are exactly the half-space $(x' - x)^\top v \ge 2\Delta$.
\end{remark}

\begin{lemma}[Closed form for adjacent pairs, Lemma 4]
  \label{lem:equality}
  \lean{MaynardZhang2026Complexity.pairValue_eq}
  \leanok
  \uses{lem:pair_value_lower, lem:pair_value_upper}
  Let $\Delta > 0$, $A$ positive definite and $(x, x') \in \I$. Then
  $p_A(x, x') = 4\Delta^2 / \norm{x - x'}^2_{A^{-1}}$.
\end{lemma}

\begin{proof}
  \leanok
  \uses{lem:pair_value_lower, lem:pair_value_upper}
  Lemmas~\ref{lem:pair_value_lower} and~\ref{lem:pair_value_upper}. (The paper states it for the
  design matrices of full rank.)
\end{proof}

\begin{lemma}[Monotonicity of the inner problem]
  \label{lem:pair_value_mono}
  \lean{MaynardZhang2026Complexity.pairValue_mono, MaynardZhang2026Complexity.pairValue_smul}
  \leanok
  \uses{def:opt_value, lem:pair_feasible_nonempty}
  For every pair $(x, x')$, $A \mapsto p_A(x, x')$ is monotone for the Loewner order on positive
  semidefinite matrices, and $p_{cA}(x, x') = c\, p_A(x, x')$ for $c \ge 0$.
\end{lemma}

\begin{proof}
  \leanok
  \uses{lem:pair_feasible_nonempty}
  Pointwise in $v$ ($A \le B$ gives $v^\top A v \le v^\top B v$); the feasible set does not depend
  on $A$, and the values are nonnegative (Lemma~\ref{lem:pair_feasible_nonempty}). (In Lean, the
  infimum over an empty feasible set is $0$ for every $A$.)
\end{proof}

\begin{lemma}[Adjacent pairs exist]
  \label{lem:exists_adjacent_pair}
  \lean{Set.mem_vertices_of_forall_norm_le, Set.Finite.exists_isAdjacent,
    MaynardZhang2026Complexity.nonempty_vertexPairs}
  \leanok
  \uses{lem:strict_expose, lem:local_global}
  A point of $\X$ of maximal norm is a vertex. If $\X$ has at least two points, it has two
  adjacent vertices; in particular $\Y \ne \emptyset$.
\end{lemma}

\begin{proof}
  \leanok
  \uses{lem:strict_expose, lem:local_global}
  If $\norm{x}$ is maximal and $y \in \X \setminus \{x\}$, then
  $0 < \norm{y - x}^2 \le 2 \norm{x}^2 - 2 y^\top x$, so $y^\top x < x^\top x$: $x$ is strictly
  exposed by itself (Lemma~\ref{lem:strict_expose}). Lemma~\ref{lem:local_global} gives a vertex
  adjacent to it.
\end{proof}

\begin{lemma}[Positivity of $f$]
  \label{lem:opt_value_pos}
  \lean{MaynardZhang2026Complexity.optValue_pos,
    MaynardZhang2026Complexity.exists_pairFeasible_lt_two_mul_optValue,
    MaynardZhang2026Complexity.bddAbove_range_iInf_pairValue,
    MaynardZhang2026Complexity.iInf_pairValue_le_optValue,
    MaynardZhang2026Complexity.le_pairValue_inv_card_smul_sum_outerSelf}
  \leanok
  \uses{def:opt_value, lem:pair_feasible_nonempty}
  If $\Delta > 0$ and $\X$ has at least two points, then $0 < f(\X, \Delta) < \infty$ (the supremum
  over designs is bounded). Consequently, for every design $\lambda$ some $(x, x') \in \Y$ has a
  feasible $(\theta, v)$ with $v^\top A(\lambda) v < 2 f(\X, \Delta)$.
\end{lemma}

\begin{proof}
  \leanok
  \uses{lem:pair_feasible_nonempty, lem:exists_adjacent_pair, lem:pair_value_lower}
  $\Y$ is nonempty (Lemma~\ref{lem:exists_adjacent_pair}) and finite. Bounded: fix a pair of $\Y$
  and a feasible $(\theta, v)$ for it; for every design,
  $p_{A(\lambda)} \le v^\top A(\lambda) v = \sum_x \lambda_x (x^\top v)^2 \le \sum_{y \in \X}
  (y^\top v)^2$. Positive: let $A_u = \frac1n \sum_{y \in \X} y y^\top$ be the design matrix of the
  uniform design, $n = \abs{\X}$. For $(x, x') \in \Y$ every feasible $v$ has
  $2\Delta \le x'^\top v - x^\top v$ (first step of Lemma~\ref{lem:pair_value_lower}), so
  $4 \Delta^2 \le 2 ((x^\top v)^2 + (x'^\top v)^2) \le 2 n\, v^\top A_u v$ (as $x, x' \in \X$), and
  $p_{A_u}(x, x') \ge 2\Delta^2 / n$; take the minimum over the finitely many pairs. Finally,
  $\min_{\Y} p_{A(\lambda)} \le f(\X, \Delta) < 2 f(\X, \Delta)$: some pair has
  $p_{A(\lambda)}(x, x') < 2 f(\X, \Delta)$, and some feasible point of that pair is below
  $2 f(\X, \Delta)$.
\end{proof}

\section{Lower bounds on the error}

\begin{lemma}[Runs agreeing before a round]
  \label{lem:actions_law_agree}
  \lean{Learning.IsAlgEnvSeq.identDistrib_action_of_env_eq,
    Learning.IsAlgEnvSeq.identDistrib_history_of_env_eq, Learning.IsAlgEnvSeqUntil.of_env_eq}
  \leanok
  Let two environments have the same feedback kernels at rounds $0, \dots, t - 1$. For a given
  algorithm, the law of $X_t$ (and of the history of the first $t$ rounds) is the same in every run
  against either environment.
\end{lemma}

\begin{proof}
  \leanok
  The law of the history of the first $t$ rounds followed by the action of round $t$ is the
  composition of the policy kernels of rounds $0, \dots, t$ and of the feedback kernels of rounds
  $0, \dots, t - 1$ (uniqueness of the law of a run).
\end{proof}

\begin{lemma}[Optimization-based lower bound, Lemma 3]
  \label{lem:opt_lower}
  \lean{MaynardZhang2026Complexity.exists_minGapGE_le_max_error_optValue,
    MaynardZhang2026Complexity.isBestArm_iff_eq_of_forall_vertices,
    MaynardZhang2026Complexity.avgParam_ite_eq, MaynardZhang2026Complexity.avgParam_add_const,
    Learning.designOfWeights, Learning.isDesign_designOfWeights}
  \leanok
  \uses{lem:gen_lower, def:opt_value, def:min_gap, lem:opt_value_pos, lem:actions_law_agree,
    lem:kl_gaussian, lem:min_gap_best_arm}
  Let $\X$ have at least two points, $\Delta > 0$, $T \ge 1$ and $A$ an identification algorithm
  with budget $T$. There are parameter sequences $\theta$, $\theta'$, both with min-gap at least
  $\Delta$, such that for every run of $A$ against $\theta$ and against $\theta'$ with standard
  Gaussian noise,
  \[
    \max\left(\Prob_\theta(\hat x \text{ is not a best arm}),
      \Prob_{\theta'}(\hat x \text{ is not a best arm})\right)
    \ge \frac14 \exp\left(-T f(\X, \Delta)\right).
  \]
\end{lemma}

\begin{proof}
  \leanok
  \uses{lem:gen_lower, lem:opt_value_pos, lem:actions_law_agree, lem:kl_gaussian,
    lem:min_gap_best_arm}
  Let $P_0$ be the law of a run of $A$ against $\theta \equiv 0$ and
  $\lambda_x = \frac1T \sum_{t < T} P_0(X_t = x)$, a design. Let $(x, x') \in \Y$ minimize
  $p_{A(\lambda)}$; since $p_{A(\lambda)}(x, x') \le f(\X, \Delta) < 2 f(\X, \Delta)$
  (Lemma~\ref{lem:opt_value_pos}), there is a feasible $(\tilde\theta, \tilde v)$ with
  $\tilde v^\top A(\lambda) \tilde v \le 2 f(\X, \Delta)$. Let
  \[
    \theta_t = 0,\ \theta'_t = \tilde v \ (t < T - 1), \qquad
    \theta_{T - 1} = T \tilde\theta,\ \theta'_{T - 1} = T \tilde\theta + \tilde v.
  \]
  Then $\thetabar_T = \tilde\theta$ and $\thetabar'_T = \tilde\theta + \tilde v$, so by feasibility
  both sequences have min-gap at least $\Delta$, with best arms $x$ and $x'$
  (Lemma~\ref{lem:min_gap_best_arm}). The law of $X_t$, $t \le T - 1$, under $\theta$ depends only
  on $\theta_0, \dots, \theta_{T - 2} = 0$, so it is its law under $P_0$
  (Lemma~\ref{lem:actions_law_agree}). The divergence between the reward laws of arm $a$ at round
  $t$ is $(a^\top \tilde v)^2 / 2$ (Lemma~\ref{lem:kl_gaussian}), so Lemma~\ref{lem:gen_lower} with
  $k = x$, $k' = x'$ gives the bound with exponent
  $\sum_{t < T} \sum_a P_0(X_t = a) (a^\top \tilde v)^2 / 2 = \frac{T}{2} \tilde v^\top A(\lambda)
  \tilde v \le T f(\X, \Delta)$. The error events $\{\hat x \ne x\}$ and $\{\hat x \ne x'\}$ are the
  events that $\hat x$ is not a best arm, the best arms being unique.
\end{proof}

\begin{remark}
  \label{rem:opt_lower_route}
  The paper takes $\theta_t = 0$, $\theta'_t = v$ in the first half of the horizon and
  $\theta_t = \theta'_t = \theta$ in the second, ``assuming without loss of generality that $T$ is
  even''. This yields the exponent $T \tilde v^\top A(\lambda) \tilde v$, so the bound
  $\exp(-T f)$ needs a feasible point attaining the inner infimum $p_{A(\lambda)}$ (true, by a
  Frank--Wolfe argument for a convex quadratic on a polyhedron, but not available in Mathlib), and
  for $T = 1$ the first half is the whole horizon, where $\thetabar = 0$ has no min-gap. The
  construction above spreads the difference between the instances over all $T$ rounds, while the
  parameters of $\theta$ before the last round vanish, so that the law of the actions under
  $\theta$ is fixed before $(\tilde\theta, \tilde v)$ is chosen. This halves the exponent, so no
  attainment is needed, and works for every $T \ge 1$.
\end{remark}

\begin{lemma}[Bound on $f$]
  \label{lem:opt_value_le}
  \lean{MaynardZhang2026Complexity.optValue_le, MaynardZhang2026Complexity.hAdjacent_pos,
    MaynardZhang2026Complexity.adjDesignValue_pos,
    MaynardZhang2026Complexity.iInf_pairValue_le_one_add_mul}
  \leanok
  \uses{lem:equality, lem:pair_value_mono, def:h_adjacent, lem:exists_adjacent_pair,
    lem:mahalanobis_cauchy_schwarz}
  If $\X$ spans $\R^d$ and has at least two points and $\Delta > 0$, then the value of the adjacent
  design is positive and $f(\X, \Delta) \le 4 / \Hadj(\X, \Delta)$.
\end{lemma}

\begin{proof}
  \leanok
  \uses{lem:equality, lem:pair_value_mono, lem:exists_adjacent_pair,
    lem:mahalanobis_cauchy_schwarz}
  Since $\X$ spans, there is a positive definite design matrix $A_0$ (the uniform design on a basis
  contained in $\X$). Positivity: let $(x, x') \in \I$ (Lemma~\ref{lem:exists_adjacent_pair}),
  $z = x - x' \ne 0$ and $K = \sum_{y \in \X} (y^\top z)^2$. For every design $\lambda$ with
  $A(\lambda)$ positive definite, Lemma~\ref{lem:mahalanobis_cauchy_schwarz} and
  $z^\top A(\lambda) z = \sum_y \lambda_y (y^\top z)^2 \le K$ give
  $\norm{z}^4 \le \norm{z}^2_{A(\lambda)^{-1}} K$, so (applied to $A_0$) $K > 0$ and the value of
  $\lambda$ is at least $\norm{z}^2_{A(\lambda)^{-1}} \ge \norm{z}^4 / K > 0$.
  Let $a$ be the value of the adjacent design, $\lambda$ a design and $\eps > 0$. Then
  $A_\eps = (A(\lambda) + \eps A_0)/(1 + \eps)$ is the positive definite design matrix of a design,
  and $A(\lambda) \le (1 + \eps) A_\eps$. Let $(x, x') \in \I$ maximize
  $\norm{x - x'}^2_{A_\eps^{-1}}$, so that $\norm{x - x'}^2_{A_\eps^{-1}} \ge a$. By
  Lemmas~\ref{lem:pair_value_mono} and~\ref{lem:equality},
  \[
    \min_{\Y} p_{A(\lambda)} \le p_{A(\lambda)}(x, x') \le (1 + \eps)\, p_{A_\eps}(x, x')
    = (1 + \eps) \frac{4 \Delta^2}{\norm{x - x'}^2_{A_\eps^{-1}}}
    \le (1 + \eps) \frac{4\Delta^2}{a}.
  \]
  Let $\eps \to 0$ and take the supremum over $\lambda$: $f(\X, \Delta) \le 4\Delta^2 / a = 4 /
  \Hadj(\X, \Delta)$.
\end{proof}

\begin{theorem}[Lower bound, Theorem 1]
  \label{thm:lower_bound}
  \lean{MaynardZhang2026Complexity.exists_minGapGE_le_max_error}
  \leanok
  \uses{lem:opt_lower, lem:opt_value_le}
  Let $\X$ span $\R^d$ and have at least two points, $\Delta > 0$, $T \ge 1$ and $A$ an
  identification algorithm with budget $T$. There are parameter sequences $\theta$, $\theta'$, both
  with min-gap at least $\Delta$, such that for every run of $A$ against $\theta$ and against
  $\theta'$ with standard Gaussian noise,
  \[
    \max\left(\Prob_\theta(\hat x \text{ is not a best arm}),
      \Prob_{\theta'}(\hat x \text{ is not a best arm})\right)
    \ge \frac14 \exp\left(-\frac{4T}{\Hadj(\X, \Delta)}\right).
  \]
\end{theorem}

\begin{proof}
  \leanok
  \uses{lem:opt_lower, lem:opt_value_le}
  Lemma~\ref{lem:opt_lower} and $T f(\X, \Delta) \le 4T / \Hadj(\X, \Delta)$
  (Lemma~\ref{lem:opt_value_le}).
\end{proof}

\begin{remark}
  \label{rem:lower_bound_degenerate}
  The paper states Theorem~\ref{thm:lower_bound} and Lemma~\ref{lem:opt_lower} without $T \ge 1$
  and without requiring two arms. Both are needed: for $T = 0$, $\thetabar_0 = 0$ and no parameter
  sequence has min-gap $\Delta > 0$; for $\X = \{x\}$ there are no adjacent pairs, so
  $\Hadj(\X, \Delta) = 0$ and $f(\X, \Delta) = 0$ (with the conventions $a / 0 = 0$ and
  $\min \emptyset = 0$), and the bound would claim an error of at least $1/4$, while every algorithm
  outputs the only arm, which is a best arm. Lemma~\ref{lem:opt_lower} holds without the assumption
  that $\X$ spans $\R^d$.
\end{remark}
